\section{问题分析}\label{sec:analysis}

\subsection{总体思路}

四个问题描述的是同一个物理过程：热量由烘房空气经表面传入药材并向内传导，水分由内部扩散到表面再进入空气。药材长度是半径的 12.5 倍，可视为无限长圆柱，温度和含水率只随半径和时间变化。因此四问共用一个一维径向传热传质模型，差别只在物性、几何和所求的量：问题一用常物性；问题二改用随含水率和温度变化的物性，并延伸到整个烘干过程；问题三在问题二的解上确定烘干完成时刻；问题四再加入半径收缩。在此基础上，本文进一步分析干燥规律、工艺参数的影响以及题设边界的能量一致性，思路见图~\ref{fig:framework}。

\begin{figure}[htbp]\centering
\includegraphics[width=0.96\textwidth]{fig_framework.pdf}
\caption{研究框架。上排为题给数据，中间为四问共用的模型与各问的侧重点，下排为在四问结果基础上的进一步分析。}
\label{fig:framework}
\end{figure}

\subsection{时间尺度与控制因素}

热量与水分由表面传到中心的特征时间分别为 $R_0^2/\alpha$ 和 $R_0^2/D$，其中 $\alpha=k/(\rho c_p)$ 为热扩散率；表面交换与内部传递的相对强弱用传热、传质 Biot 数 $Bi_h=hR_0/k$ 和 $Bi_m=h_mR_0/D$ 衡量。按附录 3 的物性估算如下。
\begin{itemize}[leftmargin=2.2em,itemsep=1pt]
\item 传热：$R_0^2/\alpha\approx\SI{0.77}{h}$，$Bi_h\approx1.0$。药材温度在约 $\SI{1}{h}$ 的量级内即可跟上烘房，表面换热与内部导热的阻力相当。
\item 传质：初始含水率下，$\SI{28}{\celsius}$ 时 $D=\SI{5.6e-9}{m^2/s}$，$\SI{50}{\celsius}$ 时 $D=\SI{1.35e-8}{m^2/s}$，$R_0^2/D\approx\SI{8.2}{h}$，$Bi_m\approx1.2$；含水率降到 $0.15$ 时 $D$ 只有 $\SI{8.0e-10}{m^2/s}$，$R_0^2/D\approx\SI{139}{h}$，$Bi_m\approx20$。
\end{itemize}
由此得到三点判断。第一，传热比传质快一到两个数量级，烘干自然分为两个阶段：先是以升温为主的预热平衡阶段，此后药材温度与烘房基本一致，进入以内部扩散为主的恒温干燥阶段。第二，含水率下降使扩散系数减小约 17 倍，传质由表面与内部共同控制逐步转为内部扩散控制，越到后期越慢，烘干时长取决于中心附近的扩散。第三，温度通过扩散系数中的因子 $\mathrm{e}^{-3850/T}$ 影响干燥快慢，由 $\SI{28}{\celsius}$ 升到 $\SI{50}{\celsius}$ 扩散系数增大 2.4 倍，烘房温度应是首要的工艺参数。附录 4 的扩散系数更小（$\SI{50}{\celsius}$、初始含水率下 $R_0^2/D\approx\SI{44}{h}$），问题四中收缩对扩散路径的缩短不可忽略。

\subsection{各问的求解要点}

\parahead{问题一} 物性为常数，扩散系数只与含水率有关，温度方程是线性的，含水率方程是非线性的，二者互不影响。$\SI{1800}{s}$ 只有热扩散时间的 0.76 倍，结果反映预热平衡阶段的前期。

\parahead{问题二} 密度、比热和导热系数随含水率变化，扩散系数随含水率和温度变化，两个方程需联立求解。除给出表 3、表 4 外，还用药材与烘房的温差识别两个阶段，用平均含水率及其下降速率描述干燥规律。

\parahead{问题三} 「各处含水率低于 0.15」等价于全域最大含水率低于 0.15。在连续时间上求最大含水率降到 0.15 的时刻作为烘干时长，结果文件输出到首个满足要求的 $\SI{60}{s}$ 采样时刻。

\parahead{问题四} 半径随时间缩小，计算域的边界随之移动。在干物质随药材均匀收缩的假设下引入参考坐标 $x=r/R(t)$，把收缩的区域映射到固定区间 $[0,1]$ 后求解，再换算到题目要求的固定位置。
